%%%%%%%%%%%%%%%%%%%%%%%%%%%%%%%%%%%%%%%%%%%%%%%%%%%%%%%%%%%%%%%%%%%%%%
% Lớp 10 | BPT và hệ  BPT bậc nhất hai ẩn | Hệ bất phương trình bậc nhất hai ẩn
% Dạng 0D2?2-1: Các khái niệm về Hệ BPT bậc nhất hai ẩn
% Số câu: 25  (sinh tự động bởi tach_cau.py; từ nay thầy sửa trực tiếp file này)
%%%%%%%%%%%%%%%%%%%%%%%%%%%%%%%%%%%%%%%%%%%%%%%%%%%%%%%%%%%%%%%%%%%%%%

%%%%%%%%%%%%%%%% TRẮC NGHIỆM 4 PHƯƠNG ÁN %%%%%%%%%%%%%%%%

% nguồn: MaHoaNH_TracNghiem_Lop10
\begin{ex}%[0D2N2-1]%[ma:57f83e8b]
	Trong các hệ bất phương trình sau, hệ bất phương trình nào là hệ bất phương trình bậc  nhất hai ẩn?
	\choice 
		{$\heva{&2024x-y > 2025\\&x-5y+z\ge  0}$}
		{$\heva{&x-y  > 0\\&x^2-x+4\le 0}$}
		{$\heva{&x+2y < 2\\&x^2-5y >  1}$}
		{\True  $\heva{&2x+y < 5\\&x-3y\ge 1}$}
	\loigiai{ 
	Theo định nghĩa hệ bất  phương trình bậc nhất hai ẩn.}
\end{ex}

% nguồn: MaHoaNH_TracNghiem_Lop10
\begin{ex}%[0D2N2-1]%[ma:841fa1c3]
	Cặp số $(x;y)=(-1;2)$ là nghiệm của hệ bất phương trình bậc nhất  hai ẩn nào dưới đây?
	 \choice
		{$\heva{&3x-y>0\\&-x+y\le 6}$} 
		{\True $\heva{&-3x-y>0\\&-x+y\le  6}$}
		 {$\heva{&-x+2y>5\\&x-y<1}$}
		{$\heva{&3x-y\le-2\\&x+y>3}$} 
	\loigiai{ 
		Cặp $(x;y)=(-1;2)$ là nghiệm của hệ bất  phương trình $\heva{&-3x-y>0\\&-x+y\le 6.}$
	} 
\end{ex}

% nguồn: MaHoaNH_TracNghiem_Lop10
\begin{ex}%[0D2N2-1]%[ma:db19c1cd]
	Hệ bất phương trình nào sau đây là hệ bất phương  trình bậc nhất hai ẩn?
	 \choice
		{$\heva{&x^2-2\le 0\\&2x+2y > 0}$} 
		{\True  $\heva{&x+y-2\le 0\\&2x-3y+2> 0}$}
		{$\heva{&3x-\dfrac{2}{y} \le 0\\&3y+2>  0}$}
		{$\heva{&x^2+y-2\le 0\\&2x-3y^2+2>  0}$}
	\loigiai{ 
		Theo định nghĩa  về hệ bất phương trình bậc nhất $2$ ẩn, ta chon đáp án $\heva{&x+y-2\le 0\\&2x-3y+2> 0.}$
	} 
\end{ex}

% nguồn: MaHoaNH_TracNghiem_Lop10
\begin{ex}%[0D2N2-1]%[ma:b8d47be4]
	Cặp số $(x;y)$ nào sau đây là nghiệm  của hệ bất phương trình
	$\heva{&x+y \le 2 \\&x - 2y \ge  4\\&x > 0.}$
	\choice 
		 {$(-1;2)$}
		{$(-2;-4)$} 
		{\True $(2;-4)$} 
		 {$(0;1)$}
	\loigiai{ 
		Thay tọa độ các điểm  vào hệ bất phương trình ta thấy điểm $(2;-4)$ thuộc miền nghiệm của cả ba bất phương trình trong hệ.
	} 
\end{ex}

% nguồn: MaHoaNH_TracNghiem_Lop10
\begin{ex}%[0D2N2-1]%[ma:789c06b4]
	Trong các cặp số sau, cặp nào \textbf{không} là nghiệm của hệ bất phương trình $\left\{\begin{array}{c}{x+y-2\le 0} \\ {2x-3y+2> 0}
\end{array}\right.$?
\choice 
	 {$\left(0; 0\right)$}
	{$\left(1;  1\right)$}
	{\True $\left(-1; 1\right)$} 
	 {$\left(-1;-1\right)$}
 \loigiai{
	Ta thay cặp số $\left(-1;  1\right)$ vào hệ ta thấy không thỏa mãn.
} 
\end{ex}

% nguồn: MaHoaNH_TracNghiem_Lop10
\begin{ex}%[0D2N2-1]%[ma:83174422]
	Cặp số nào dưới đây \textbf{không} là nghiệm  của hệ bất phương trình $\heva{&x+2y > 3\\
	&2x-y\le  1}$?
	\choice 
		{$\left(0;2\right)$} 
		{$\left(1;2\right)$} 
		 {\True $\left(1;1\right)$}
		{$\left(-1;3\right)$} 
	\loigiai{ 
		Thay lần lượt các cặp số trong các đáp án vào hệ ta thấy chỉ có cặp số $\left(1;1\right)$ không  thỏa mãn hệ.
	} 
\end{ex}

% nguồn: MaHoaNH_TracNghiem_Lop10
\begin{ex}%[0D2N2-1]%[ma:5c44e888]
	Cặp  số $(x;y)$ nào sau đây không là nghiệm của hệ bất phương trình
	$\heva{&x+y \le 2 \\&x \ge -3.}$ 
	 \choice
		{\True  $(3;0)$}
		 {$(-1;1)$}
		{$(-2;0)$} 
		{$(0;0)$} 
	\loigiai{ 
		Xét  hệ bất phương trình
		$$\heva{x+y  \le 2 \\
		x  \ge -3 \\
		y  \ge -1}$$
		Với $x=3$; $y=0$ không là nghiệm của bất phương trình thứ nhất nên $(3;0)$  không là nghiệm của hệ.
	 }
\end{ex}

% nguồn: MaHoaNH_TracNghiem_Lop10
\begin{ex}%[0D2N2-1]%[ma:d34e49e6]
	 Trong các hệ bất phương trình sau, hệ bất phương trình nào là hệ bất phương trình bậc nhất hai ẩn?
	\choice 
		 { $\heva{& 4x+5y<-10 \\
		 & {9^3}x+5y^2\ge 6}$ }
		{ $\heva{& x+y\le 1  \\
		& x+3y+z>2 \\ 
		&  x-z>0} $}
		{ $\heva{& x>3 \\ 
		& y\le 1  \\
		& x-3y^2\ge  0}$}
		{\True $\heva{& 2x-3y<0 \\ 
		 & 5x+3y+1\le 0}$}
	 \loigiai{
	Hệ bất phương trình  bậc nhất $2$ ẩn là $1$ hệ gồm $2$ hay nhiều bất phương trình bậc nhất $2$ ẩn.}
\end{ex}

% nguồn: MaHoaNH_TracNghiem_Lop10
\begin{ex}%[0D2H2-1]%[ma:65daed1c]
	Miền nghiệm của hệ bất phương  trình $\heva{&3x+y\ge 6\\&x\ge y-3}$ là phần mặt phẳng chứa điểm
	\choice 
		{\True  $(2;1)$}
		 {$(1;2\,024)$}
		 {$(0;0)$}
		 {$(1;1)$}
	 \loigiai{
	Thế tọa độ các điểm vào từng bất phương trình trong hệ ta thấy đáp án A có  toạ độ thoả mãn các bất phương trình trong hệ.}
\end{ex}

% nguồn: MaHoaNH_TracNghiem_Lop10
\begin{ex}%[0D2H2-1]%[ma:ca2f7575]
	Điểm $O(0;0)$ \textbf{không} thuộc  miền nghiệm của hệ bất phương trình nào sau đây?
	\choice 
		{\True  $\heva{&x+3y<0\\&2x+y+4>0}$}
		{$\heva{&x+3y\ge  0\\&2x+y-4<0}$}
		 {$\heva{&x+3y-6<0\\&2x+y+4>0}$}
		{$\heva{&x+3y-6<0\\&2x+y+4\ge 0}$} 
	\loigiai{Thay tọa  độ $O$ vào $\heva{&x+3y<0\\&2x+y+4>0}$ ta được $\heva{&0<0\\&04>0}$ nên $O(0;0)$ không thuộc miền nghiệm của hệ bất phương trình này.
	} 
\end{ex}

% nguồn: MaHoaNH_TracNghiem_Lop10
\begin{ex}%[0D2H2-1]%[ma:c83d4dbc]
	 Điểm nào sau đây thuộc miền nghiệm của hệ bất phương trình $\heva{&2x-y > 0\\&x-3y < -3\\&x+y > 6}$?
	 \choice
		{$A(1;1)$} 
		 {\True $B(7;4)$}
		 {$C(2;1)$}
		{$D(-2;4)$} 
	 \loigiai{
	Toạ  độ điểm $B$ thảo mãn hệ bất phương trình đã cho.}
\end{ex}

% nguồn: MaHoaNH_TracNghiem_Lop10
\begin{ex}%[0D2H2-1]%[ma:0e26213c]
	Cặp số $(x;y)$ nào sau đây không là nghiệm của hệ bất phương trình $\heva{&3x+4y-1> 0 \\&x+2y-3\le 0.}$ 
	 \choice
		 {$(3;0)$}
		{$(-1;2)$} 
		 {$(2;0)$}
		 {\True $(0;0)$}
	\loigiai{Xét hệ bất phương trình $\heva{&3x+4y-1> 0\quad(1) \\&x+2y-3\le 0\quad(2).}$\\ 
		\begin{itemize}
			\item Với $x=3$; $y=0$ là nghiệm của bất phương trình $(1)$ và $(2)$ nên $(3;0)$ là nghiệm của  hệ.
			\item Với $x=-1$; $y=2$ là nghiệm của bất phương trình $(1)$ và $(2)$ nên $(-1;2)$ là nghiệm của hệ. 
			\item Với $x=2$; $y=0$ là nghiệm của bất phương trình $(1)$  và $(2)$ nên $(2;0)$ là nghiệm của hệ.
			\item Với $x=0$; $y=0$ không là nghiệm của bất phương trình $(1)$ nên $(0;0)$ không là nghiệm  của hệ.
	\end{itemize}}
\end{ex}

% nguồn: MaHoaNH_TracNghiem_Lop10
\begin{ex}%[0D2H2-1]%[ma:a211b50b]
	 Cặp số $\left(x;y\right)$ nào sau đây là nghiệm của hệ bất phương trình $\heva{&x+3y\le 0\\
	 &x-2y\ge 1.}$
	 \choice
		 {$\left(1;0\right)$}
		 {$\left(3;1\right)$}
		{\True $\left(1;-3\right)$} 
		{$\left(-2;3\right)$} 
	\loigiai{Xét hệ bất phương trình  $\heva{&x+3y\le 0 \quad(1)\\&x-2y\ge 1. \quad (2)}$\\
		\begin{itemize}
			\item Với $x=1$; $y=0$ không là nghiệm của bất phương trình $(1)$ nên  $(1;0)$ không là nghiệm của hệ.
			\item Với $x=3$; $y=1$ không  là nghiệm của bất phương trình $(1)$ nên $(3;1)$ không là nghiệm của hệ.
			\item Với $x=1$; $y=-3$ là nghiệm của cả hai bất phương trình $(1)$ và $(2)$ nên $(1;-3)$ là  nghiệm của hệ.
			\item Với $x=-2$; $y=3$ không là nghiệm của bất phương trình  $(1)$ nên $(-2;3)$ không là nghiệm của hệ.
	\end{itemize}}
\end{ex}

% nguồn: MaHoaNH_TracNghiem_Lop10
\begin{ex}%[0D2B2-1]%[ma:88447d8c]
	 Có bao nhiêu cách sắp xếp $4$ học sinh thành một hàng dọc?
	\choice 
		{$\mathrm{C}_{4}^{1}$} 
		{$16$} 
		{$8$} 
		 {\True $24$}
	 \loigiai{
		 Mỗi cách xếp thoả đề chính là một hoán vị của $4$ học sinh. Vậy có $4!=24$ cách xếp.
	} 
\end{ex}

% nguồn: MaHoaNH_TracNghiem_Lop10
\begin{ex}%[0D2Y2-1]%[ma:6956c41c]
	Có bao nhiêu  cách xếp $4$ học sinh vào bốn chiếc ghế được xếp thành hàng ngang?
	\choice 
		 {\True $4!$}
		 {$4^4$}
		{$16$} 
		 {$10$}
	\loigiai{ 
		Xếp $4$ học sinh vào bốn  chiếc ghế có $4!$ cách.
	 }
\end{ex}

% nguồn: MaHoaNH_TracNghiem_Lop10
\begin{ex}%[0D2Y2-1]%[ma:438e1b64]
	Số cách sắp xếp $6$ học sinh  ngồi vào một dãy gồm $6$ ghế là
	 \choice
		{\True  $6!$}
		 {$\mathrm{A}_6^1$}
		 {$\mathrm{C}_6^6$}
		{$6^6$} 
	 \loigiai{
		Số cách sắp xếp $6$ học  sinh ngồi vào một dãy gồm $6$ ghế là $6!$.
	} 
\end{ex}


%%%%%%%%%%%%%%%% ĐÚNG/SAI %%%%%%%%%%%%%%%%

% nguồn: MaHoaNH_DungSai_Lop10
\begin{ex}%[0D2N2-1]%[ma:d3a23320]
	Cho hệ bất phương trình $\heva{&x+y<2\\&2x-4y>5}$ (I). 
	 \choiceTF
		{ Cặp số $(0;-1)$  là nghiệm của (I)}
		{  Cặp số $(2;0)$ là nghiệm của (I)}
		{ \True Cặp số $(1;-1)$  là nghiệm của (I)}
		{Cặp số $(1;-1)$ là nghiệm của  (I)
		} 
	 \loigiai{
		\begin{itemchoice}
			\itemch Cặp số $(0;-1)$  không là nghiệm của (I).
			\itemch Cặp số $(2;0)$ không là nghiệm của (I). 
			\itemch Cặp số  $(1;-1)$ là nghiệm của (I).
			\itemch Cặp số  $(1;-1)$ không là nghiệm của (I)
		\end{itemchoice}
	 }
\end{ex}

% nguồn: MaHoaNH_DungSai_Lop10
\begin{ex}%[0D2H2-1]%[ma:3d502e3a]
	Xét tính đúng, sai của các mệnh  đề sau
	 \choiceTF
		{$\heva{&x+\sqrt{y}  > 3 \\ &x^{2}-y \ge -2}$ là hệ bất phương trình bậc nhất hai ẩn}
		{\True $\heva{&x > 5 \\ &y <-2 \\ & x+y \ge 100}$  là hệ bất phương trình bậc nhất hai ẩn}
		{\True $\heva{&3^{5}x-\sqrt{20}y > 7 \\ &x+y \le  100}$ là hệ bất phương trình bậc nhất hai ẩn}
		{$\heva{&x+y+z < 10 \\  &x+y < 5 \\& 2x+3y \ge 20}$ là hệ bất phương trình bậc nhất hai ẩn}
	 \loigiai{
		\begin{itemchoice}
			\itemch Vì $\heva{&x+\sqrt{y} > 3 \\ &x^{2}-y \ge -2}$ chứa $\sqrt{x}$, $x^2$ nên không phải là  hệ bất phương trình bậc nhất hai ẩn.
			\itemch Vì $\heva{&x > 5  \\ &y <-2 \\ & x+y \ge 100}$ là hệ bất phương trình bậc nhất hai ẩn.
			\itemch Vì $\heva{&3^{5}x-\sqrt{20}y > 7 \\ &x+y \le  100}$ là hệ bất phương trình bậc nhất hai ẩn
			\itemch Vì $\heva{&x+y+z < 10 \\ &x+y < 5 \\& 2x+3y \ge 20}$ chứa $z$ nên là hệ bất phương  trình bậc nhất hai ẩn.
	\end{itemchoice}.}
\end{ex}

% nguồn: MaHoaNH_DungSai_Lop10
\begin{ex}%[0D2H2-1]%[ma:b083e738]
	Xét tính đúng, sai của các mệnh đề sau 
	\def\dotEX{} 
	\choiceTF 
		{Điểm $M(1;2)$ không thuộc miền nghiệm của hệ bất phương  trình $
		\heva{&x+3y-6 > 0  \\&2x+y+4 > 0.}$}
		{\True Điểm $M(1;2)$ không thuộc miền  nghiệm của hệ bất phương trình $
		\heva{&x+3y-6  < 0 \\&2x+y-4 < 0.}$}
		 {\True Điểm $(0;0)$ là một nghiệm của hệ bất phương trình $
		 \heva{& x+y-2 \le 0 \\& 2x-3y+2 > 0.}$}
		{\True Điểm $(1;1)$ là một  nghiệm của hệ bất phương trình $
		\heva{&x+y-2 \le 0 \\&2x-3y+2  > 0.}$}
	 \loigiai{
		\begin{itemchoice}
			\itemch Thay $x=1$, $y=2$ vào hệ bất phương trình, ta được $\heva{&1+3\cdot 2-6>0\\&2\cdot 1+2+4>0}$ (đúng) nên suy ra điểm $M(1;2)$ thuộc  miền nghiệm của hệ đã cho.
			\itemch Thay $x=1$, $y=2$ vào hệ bất phương trình, ta được  $\heva{&1+3\cdot 2-6<0\\&2\cdot 1+2-4<0}$ (sai) nên suy ra điểm $M(1;2)$ không thuộc miền nghiệm của hệ đã cho.
			\itemch Thay $x=0$, $y=0$ vào hệ  bất phương trình, ta được $\heva{&0+0-2\le 0\\&2\cdot 0-3\cdot 0+2>0}$ (đúng) nên suy ra điểm $(0;0)$ là một nghiệm của hệ đã cho.
			\itemch Thay $x=1$, $y=1$ vào hệ bất  phương trình, ta được $\heva{&1+1-2\le 0\\&2\cdot 1-3\cdot 1+2>0}$ (đúng) nên suy ra điểm $(1;1)$ là một nghiệm của hệ đã cho.
		\end{itemchoice}
	 }
\end{ex}

% nguồn: MaHoaNH_DungSai_Lop10
\begin{ex}%[0D2H2-1]%[ma:cede3bcb]
	Xét tính đúng, sai của các mệnh  đề sau:
	\def\dotEX{} 
	 \choiceTF
		{\True $(3;-1)$ là một nghiệm của  hệ bất phương trình $
		\heva{& x+5y < 5  \\&-3x-y \le -7.}$}
		{\True $(3;-1)$  không là một nghiệm của hệ bất phương trình $
		\heva{& x > 5 \\&y  < 4 \\&x+y \le 10.}$}
		{\True $(3;-1)$  không là một nghiệm của hệ bất phương trình $
		\heva{& -x+5y > 1 \\&3x+y  > 5.}$}
		{$(0;0)$ là một một nghiệm của hệ bất phương trình $\heva{&2x+y > 3 \\&  -x+3y \le 5 \\& 3x-y \ge 7.}$}
	\loigiai{ 
		\begin{itemchoice}
			\itemch Thay $(3;-1)$ vào hệ bất phương trình $\heva{&x+5y < 5 \\&-3x-y \le -7}$ ta được $\heva{&3+5\cdot (-1)<5\\&-3\cdot 3-(-1)\le  -7}$ (đúng).\\
			Vậy $(3;-1)$ là một  nghiệm của hệ bất phương trình.
			\itemch Thay $(3;-1)$ vào hệ  bất phương trình $
			\heva{&x > 5 \\&y < 4  \\&x+y \le 10}$ ta được
			$\heva{&3>5\\&-1<4\\&3-1\le  10}$ (sai).\\
			Vậy $(3;-1)$ không  là nghiệm của hệ bất phương trình.
			\itemch Miền nghiện của hệ bất phương trình $\heva{&-x+5y > 1 \\&3x+y > 5}$ là miền không bị gạch ở hình sau  (không kể bờ $d_1$, $d_2$)
			\begin{center}
				\begin{tikzpicture}[line join=round,line cap=round,>=stealth,thick]
						\tikzset{every  node/.style={scale=0.9}}
						\begin{scope}
						\clip (-2,-2)  rectangle (5,5);
						 \fill[pattern=north east lines] (-12,-2.2)--(25,-2.2)--(25,5.2)--cycle;
						\fill[pattern=north east  lines] (-3,14)--(-3,-13)--(6,-13)--cycle;
						\draw (24,5)--(-11,-2) node [pos=0.45,above,sloped] {$-x+5y-1=0$}; 
						\draw (0,5)--(2.33,-2) node  [pos=0.45,above,sloped] {$3x+y-5=0$};
						\end{scope}
						\draw[->] (-2,0)--(5,0)  node[below]{$x$};
						\draw[->]  (0,-2)--(0,5) node[left]{$y$};
						\fill (0,0) circle (1pt) node[below  left]{$O$};
					\end{tikzpicture}
			\end{center}
			\itemch Miền nghiệm của hệ bất phương trình $\heva{&2x+y > 3 \\&-x+3y \le 5 \\&3x-y \ge 7}$ là miền không  bị gạch ở hình sau (không kể bờ $d_1$)
			\begin{center}
				\begin{tikzpicture}[line join=round,line cap=round,>=stealth]
						\tikzset{every node/.style={scale=0.9}} 
						\begin{scope}
						\clip  (-1,-3) rectangle (5,5);
						\fill[pattern=north east lines] (-2,7)--(-2,-9)--(6,-9)--cycle; 
						\fill[pattern=north  east lines] (-21,-5.33)--(-21,5.33)--(11,5.33)--cycle;
						\fill[pattern=north east lines]  (-2,-13)--(-2,11)--(6,11)--cycle;
						\draw (-1,5)--(4,-5) node [pos=0.35,above,sloped] {$2x+y-3=0$}; 
						\draw (10,5)--(-20,-5) node [pos=0.25,above,sloped] {$-x+3y-5=0$}; 
						\draw (4,5)--(0.67,-5) node [pos=0.55,above,sloped] {$3x-y-7=0$}; 
						\end{scope}
						 \draw[->] (-1,0)--(5,0) node[below]{$x$};
						\draw[->] (0,-3)--(0,5)  node[left]{$y$};
						 \fill (0,0) circle (1pt) node[below left]{$O$};
					\end{tikzpicture}
			\end{center}
		\end{itemchoice}
	 }
\end{ex}


%%%%%%%%%%%%%%%% TRẢ LỜI NGẮN %%%%%%%%%%%%%%%%

% nguồn: MaHoaNH_TraLoiNgan_Lop10
\begin{bt}%[0D2H2-1]%[ma:1d8fdb77]
	Cho các cặp số sau: $(0; 0)$; $(1; 1)$; $(-1; 1)$; $(-1;-1)$. Hỏi có bao nhiêu  cặp không là nghiệm của hệ bất phương trình $\heva{&x+y-2 \leq 0 \\& 2 x-3 y+2>0.}$
	\shortans{1} 
	\loigiai{Ta thay cặp số $(-1; 1)$ vào hệ ta  thấy không thỏa mãn.
		Vậy chỉ có một cặp số không là nghiệm  của hệ bất phương trình.
	} 
\end{bt}

% nguồn: MaHoaNH_TraLoiNgan_Lop10
\begin{bt}%[0D2V2-1]%[ma:35aaeaf3]
	Tìm GTLN của $f\left(x,  y\right)=x+2 y$ với điều kiện
	$$\heva{&0 \le y \le 4 \quad\quad\quad \left(d_{1}\right)\\&0 \le x \quad\quad\quad\quad\quad \left(d_{2}\right)\\ &x-y-1 \le  0 \quad\quad \left(d_{3}\right)\\ &x+2y-10 \le 0 \quad \left(d_{4}\right).}$$
	\shortans{10} 
	\loigiai{ 
		\begin{center}
			\begin{tikzpicture}[scale=0.7,line join=round,line cap=round,>=stealth]
					\tikzset{every node/.style={scale=1}} 
					\def\xmin{-3}\def\xmax{7}\def\ymin{-1}\def\ymax{6} 
					\draw[->] (\xmin-0.2,0)--(\xmax+0.2,0) node[below]{$x$}; 
					 \draw[->] (0,\ymin-0.2)--(0,\ymax+0.2) node[right]{$y$};
					 \draw (0,0) node[below right]{$O$};
					\clip (\xmin,\ymin) rectangle (\xmax,\ymax); 
					\draw[thick,smooth,samples=200,domain=\xmin:\xmax] plot  (\x,{0*(\x)+4});
					\draw[thick,smooth,samples=200,domain=\xmin:\xmax] plot  (\x,{1*(\x)-1});
					\draw[thick,smooth,samples=200,domain=\xmin:\xmax] plot  (\x,{(-1/2)*(\x)+5});
					\fill[pattern=north east  lines,smooth]
						plot[domain=-3:7]  (\x,{0*(\x)+6})
						-- plot[domain=7:-3] (\x,{0*(\x)+4}) 
						-- cycle; 
					\fill[pattern=north  east lines,smooth]
						plot[domain=-2:7]  (\x,{0*(\x)+6})
						 -- plot[domain=7:-2] (\x,{(-1/2)*(\x)+5})
						 -- cycle;
					\fill[pattern=north  east lines,smooth]
						 plot[domain=0:7] (\x,{1*(\x)-1})
						-- plot[domain=7:0] (\x,{0*(\x)-2}) 
						-- cycle; 
					\fill[pattern=north east  lines,smooth]
						 plot[domain=-5:0] (\x,{0*(\x)+6})
						-- plot[domain=0:-5]  (\x,{0*(\x)-2})
						-- cycle; 
					\fill[pattern=north  east lines,smooth]
						 plot[domain=-5:7] (\x,{0*(\x)})
						--  plot[domain=7:-5] (\x,{0*(\x)-2})
						--  cycle;
					\path 
						(0,4) node[below right]{$D$} 
						(1,0) node[above  left]{$A$}
						(4,3)  node[left]{$B$}
						 (2,4) node[below]{$C$};
				\end{tikzpicture}
		\end{center}
		$M\left(x; y\right)$ thoả  mãn (I) là miền bên trong đa giác $OABCD$.\\
		Tìm toạ độ $A$, $B$, $C$, $D$ bằng phương pháp đồ thị hay phương trình hoành độ  giao điểm.\\
		 Thay toạ độ $O$, $A$, $B$, $C$, $D$ vào $f\left(x; y\right)=x+2 y$ ta có
		\begin{center}
			\begin{tabular}{|c|c|c|c|c|c|}
				\hline & $O$ & $A$ & $B$ &  $C$ & $D$ \\
				\hline$f(x; y)$  & $0$ & $1$ & $10$ & $10$ & $8$ \\
				 \hline
			\end{tabular}
		\end{center}
		Vậy $\max f\left(x\right)=10$. 
	 }
\end{bt}


%%%%%%%%%%%%%%%% TỰ LUẬN %%%%%%%%%%%%%%%%

% nguồn: MaHoaNH_TuLuan_Lop10
\begin{bt}%[0D2H2-1]%[ma:91b83d1e]
	Hệ bất phương trình nào sau đây là hệ bất phương trình bậc nhất  hai ẩn?
	\begin{enumEX}[a)]{3}
		\item $\heva{&x < 0 \\&y\ge  0}$;
		\item $\heva{&x+y^2< 0 \\&y-x > 1}$; 
		\item $\heva{&x+y+z <  0 \\&y < 0}$;
		\item $\heva{&-2x+y < 3^2 \\&4^2x+3y  < 1}$;
		\item $\heva{&x+2y\ge-1 \\&2x^2-2y\le-3}$; 
		\item  $\heva{&-x+3y\ge-2 \\& 5\left|x\right|+3y^2\le 4.}$
	\end{enumEX}
	 \loigiai{
		\begin{enumEX}[a)]{1}
			 \item Ta thấy hệ $\heva{&x < 0 \\&y\ge 0}$ gồm hai bất phương trình bậc nhất hai ẩn là $x < 0$ và $y\ge 0$ nên hệ trên là hệ bất phương trình bậc nhất hai ẩn.
			\item Ta thấy hệ $\heva{&x+y^2< 0 \\ &y-x > 1}$ không là hệ bất phương trình bậc nhất hai ẩn vì $x+y^2< 0$ không là bất phương trình  bậc nhất hai ẩn (chứa $y^2$).
			\item Ta thấy hệ $\heva{&x+y+z < 0 \\ &y < 0}$ không là hệ bất phương  trình bậc nhất hai ẩn vì $x+y+z < 0$ không là bất phương trình bậc nhất hai ẩn (có $3$ ẩn).
			\item Ta có $\heva{&-2x+y < 3^2\\&4^2x+3y < 1}\Leftrightarrow \heva{&-2x+y < 9 \\&16x+3y  < 1}$ đây là hệ bất phương trình bậc nhất hai ẩn và gồm hai bất phương trình bậc nhất hai ẩn là $-2x+y < 9$ và $16x+3y < 1$.
			\item Do bất phương trình $2x^2-2y\le-3$ không là bất phương trình bậc nhất hai  ẩn, nên hệ bất phương trình đã cho không là hệ bất phương trình bậc nhất hai ẩn.
			\item Do bất phương trình $5\left|x\right|+3y^2\le 4$ không là bất phương trình bậc nhất hai ẩn, nên hệ bất phương trình  đã cho không là hệ bất phương trình bậc nhất hai ẩn.
	\end{enumEX}}
\end{bt}

% nguồn: MaHoaNH_TuLuan_Lop10
\begin{bt}%[0D2G2-1]%[ma:2ec82248]
	Xếp ngẫu nhiên $2$ học sinh lớp 10A, $3$ học sinh lớp 10B và $5$ học sinh  lớp 10C thành một hàng ngang. Hỏi có bao nhiêu cách xếp để không có học sinh nào của cùng một lớp đứng cạnh nhau.
	\loigiai{ 
		Số phần tử  của không gian mẫu là $n(\Omega)=10!$.\\
		 Gọi $A$ là biến cố \lq\lq Không có học sinh nào của cùng một lớp đứng cạnh nhau\rq\rq.\\
		Sắp xếp $5$ học sinh lớp 10C  có $5!$ cách.\\
		Ứng với mỗi cách sắp xếp $5$ học sinh lớp 10C sẽ có $6$ khoảng trống gồm $4$ vị trí ở giữa và hai vị trí ở hai đầu để xếp các học sinh còn  lại.
		\begin{enumerate}
			\item Trường hợp 1. Xếp $3$ học sinh 10B vào $4$ vị trí ở giữa  (không xếp vào hai đầu), có $\mathrm{A}_4^3$ cách.\\
			Ứng với mỗi cách xếp đó, chọn $1$ trong $2$ học sinh  lớp 10A xếp vào $4$ vị trí trống để hai học sinh 10C không được cạnh nhau có $2$ cách.\\
			Học sinh  10A còn lại còn có $8$ vị trí để sắp xếp.\\
			Theo quy tắc nhân ta có $5!\cdot  \mathrm{A}_4^3\cdot 2\cdot 8$ cách.
			\item Xếp  $2$ trong $3$ học sinh lớp 10B vào $4$ vị trí trống ở giữa và học sinh còn lịa xếp vào $2$ cầu có $\mathrm{A}_4^2\cdot 2\cdot \mathrm{C}_2^1$.\\
			Ứng với mỗi cách xếp đó sẽ còn $2$ vị trí ở giữa, xếp hai học sinh lớp 10A vào  đó có $2!$ cách.\\
			 Theo quy tắc nhân, có $\mathrm{A}_4^2\cdot 2\cdot \mathrm{C}_2^1\cdot 2!$ cách.
		\end{enumerate}
		 Suy ra $n(A)=5!\cdot \mathrm{A}_4^3\cdot 2\cdot 8+\mathrm{A}_4^2\cdot 2\cdot \mathrm{C}_2^1\cdot 2!=63360$ cách.\\
		Vậy $\mathrm{P}(A)=\dfrac{n(A)}{n(\Omega)}=\dfrac{63360}{10!}=\dfrac{11}{630}$. 
	 }
\end{bt}

% nguồn: MaHoaNH_TuLuan_Lop10
\begin{bt}%[0D2B2-1]%[ma:d6c8210a]
	Biểu diễn miền  nghiệm của hệ bất phươmg trình:
	$\heva{&3x+y \leq 6 \\&x+y \leq  4 \\&x \geq 0 \\&y \geq 0.}$
	\loigiai{ 
	\immini{Biểu diễn từng miền nghiệm  của mỗi bất phương trình trên mặt phẳng $Oxy$.\\
		Miền không gạch chéo (miền tứ giác $OABC$, bao gồm cả các cạnh) trong hình bên là phần giao của các miền nghiệm và cũng là phần biểu diễn miền nghiệm của hệ bất phương trình  đã cho.
		 }{
		\begin{tikzpicture}[line join=round,line cap=round,>=stealth,thick]
				\tikzset{label  style/.style={font=\footnotesize}}
				\begin{scope}
					\clip (-3,-1) rectangle  (5,7);
					 \fill[pattern=north east lines] (-4,18)--(6,18)--(6,-12)--cycle;
					\fill[pattern=north  east lines] (-4,8)--(6,8)--(6,-2)--cycle;
					\fill[pattern=north east  lines] (0,-1)--(-3,-1)--(-3,7)--(0,7)--cycle;
					\fill[pattern=north east lines] (-3,0)--(-3,-1)--(5,-1)--(5,0)--cycle; 
					\draw (-0.33,7)--(2.33,-1) node  [pos=0.3,above,sloped] {$3x+y-6=0$};
					 \draw (-3,7)--(5,-1) node [pos=0.7,above,sloped] {$x+y-4=0$};
				\end{scope}
				\draw[->]  (-3,0)--(5,0) node[below]{$x$};
				 \draw[->] (0,-1)--(0,7) node[left]{$y$};
				\draw (0,0) node[below left]{$O$} 
					(0,4) node[right]{$A$} 
					 (1,3) node[right]{$B$}
					 (2,0) node[above right]{$C$};
				 \draw[dashed](1,0)--(1,3)--(0,3);
				\foreach \x  in {-2,-1,1,2,3,4}
				\draw[thin] (\x,1pt)--(\x,-1pt) node [below] {$\x$}; 
				\foreach \y  in {1,2,3,4,5,6}
				\draw[thin] (1pt,\y)--(-1pt,\y) node [left] {$\y$}; 
			\end{tikzpicture}
	 }}
\end{bt}
